% ============================================================================
%  4. DESARROLLO T\'ECNICO DEL PROYECTO
%  ----------------------------------------------------------------------------
%  Integraci\'on del contenido de las actividades a4 (correo y archivos),
%  a5 (DHCP y Web) y a7 (Samba 4 AD DC).
% ============================================================================

\section{Desarrollo t\'ecnico del proyecto}

Este cap\'itulo documenta la implementaci\'on de los servicios de
infraestructura del banco Sudoers. El desarrollo se organiz\'o en tres
bloques correspondientes a las actividades del curso: primero el servicio de
correo electr\'onico y el de archivos multi-protocolo, luego el servidor DHCP
y el servidor Web con proxy inverso, y finalmente el despliegue del
controlador de dominio Samba 4 Active Directory que integra la identidad de
toda la red.

Cada bloque describe desde la instalaci\'on de los paquetes necesarios hasta
las pruebas funcionales ejecutadas desde los clientes de la LAN. Las
direcciones IP citadas corresponden al plan de direccionamiento final
documentado en el cap\'itulo de arquitectura: el host \texttt{dc1} opera en
\texttt{192.168.0.2}, el servidor de monitoreo \texttt{zabbix} en
\texttt{192.168.0.3} y el router en \texttt{192.168.0.1}.

% ============================================================================
%  BLOQUE 1: SERVIDOR DE CORREO ELECTR\'ONICO
% ============================================================================
\section{Servidor de correo electr\'onico}

El servicio de correo se implement\'o sobre Debian 12 mediante dos
componentes instalados de forma nativa: Postfix como agente de
transferencia de correo (\textsc{mta}) y Dovecot como agente de acceso a
buzones (\textsc{mda}). La protecci\'on de las conexiones se realiza con
\textsc{starttls} y certificados auto-firmados \cite{postel1982}.

\subsection{Arquitectura del servicio}

Postfix gestiona el env\'io y la recepci\'on mediante SMTP, habilitando el
puerto 25 para la transferencia entre servidores y el 587 para el env\'io
autenticado desde clientes. Dovecot expone los buzones mediante IMAP
(puerto 143) y POP3 (puerto 110), y valida las credenciales de los usuarios
a trav\'es del subsistema SASL.

La autenticaci\'on se delega de Postfix a Dovecot mediante un socket Unix
(\texttt{/var/run/dovecot/private/auth}). El cliente remite sus
credenciales en el puerto 587, Postfix consulta al socket de Dovecot y s\'olo
tras la validaci\'on exitosa acepta el mensaje. La figura
\ref{fig:arquitectura-correo} muestra el esquema del servicio.

\figuraAPA{fig:arquitectura-correo}{Arquitectura del servicio de correo: Postfix, Dovecot y DNS}{%
  \begin{tikzpicture}[
      node distance=1.2cm and 1.8cm,
      servicio/.style={draw, rounded corners, minimum width=2.6cm,
                       minimum height=0.9cm, align=center, font=\small},
      cliente/.style={draw, circle, minimum size=1cm, font=\small},
      linea/.style={-Stealth, thick}
    ]
    \node[servicio] (postfix) {Postfix\\{\footnotesize SMTP}};
    \node[servicio, right=of postfix] (dovecot) {Dovecot\\{\footnotesize IMAP/POP3}};
    \node[cliente, left=of postfix] (c1) {C1};
    \node[cliente, below=of c1] (c2) {C2};
    \node[servicio, below=of postfix, yshift=-0.4cm] (dns) {BIND9\\{\footnotesize DNS}};
    \draw[linea] (c1.east) -- node[above,font=\scriptsize]{587} (postfix.west);
    \draw[linea] (c2.east) -| (postfix.west);
    \draw[linea] (postfix.east) -- node[above,font=\scriptsize]{SASL} (dovecot.west);
    \draw[linea, dashed] (postfix.south) -- (dns.north);
  \end{tikzpicture}%
}{Elaboraci\'on propia.}

\subsection{Instalaci\'on de paquetes}

\begin{lstlisting}[style=terminal]
$ apt update
$ apt install -y postfix dovecot-core dovecot-lmtpd \
  libsasl2-modules openssl
\end{lstlisting}

El paquete \texttt{postfix} aporta el daemon \texttt{master} y el gestor
\texttt{postconf}; \texttt{dovecot-core} y \texttt{dovecot-lmtpd} instalan el
servidor de buzones con soporte \textsc{lmtp}; \texttt{libsasl2-modules}
habilita los mecanismos de autenticaci\'on; y \texttt{openssl} se emplea
para emitir los certificados \textsc{tls} \cite{opensslDocs}.

\subsection{Configuraci\'on del servidor DNS}

Antes de configurar el correo se estableci\'o la resoluci\'on de nombres del
dominio \texttt{sudoers.lan}, con un registro \textsc{mx} que apunta al
servidor de correo y su registro \textsc{a} correspondiente. La
verificaci\'on se realiz\'o con la utilidad \texttt{dig}:

\begin{lstlisting}[style=terminal]
$ dig +short MX sudoers.lan
10 mail.sudoers.lan.

$ dig +short mail.sudoers.lan
192.168.0.2
\end{lstlisting}

\capturaAPA{fig:dig-mx}{Verificaci\'on del registro \textsc{mx} del dominio \texttt{sudoers.lan} con \texttt{dig}}{captura-dig-mx.png}{Captura de pantalla de la salida de \texttt{dig +short MX sudoers.lan} y \texttt{dig +short mail.sudoers.lan} en la terminal, mostrando el registro MX y la direcci\'on IP del servidor de correo.}

\subsection{Configuraci\'on del servidor SMTP}

La configuraci\'on de Postfix reside en \texttt{/etc/postfix/main.cf}. Los
par\'ametros relevantes se resumen en la tabla \ref{tab:parametros-postfix}.

\begin{table}[H]
\centering
\caption{Par\'ametros principales de Postfix}
\label{tab:parametros-postfix}
\contenidoConFuente{%
  \begin{tablaAPA}
    \begin{tabular}{@{}p{5.8cm}p{8cm}@{}}
    \toprule
    \textbf{Par\'ametro} & \textbf{Valor} \\
    \midrule
    \texttt{myhostname} & \texttt{mail.sudoers.lan} \\
    \texttt{mydomain} & \texttt{sudoers.lan} \\
    \texttt{mynetworks} & \texttt{127.0.0.0/8, 10.0.0.0/8,}\newline
                          \texttt{172.16.0.0/12, 192.168.0.0/16} \\
    \texttt{home\_mailbox} & \texttt{Maildir/} \\
    \texttt{smtpd\_sasl\_type} & \texttt{dovecot} \\
    \texttt{smtpd\_sasl\_path} & \texttt{private/auth} \\
    \texttt{smtpd\_sasl\_auth\_enable} & \texttt{yes} \\
    \texttt{smtpd\_sasl\_security\_options} & \texttt{noanonymous} \\
    \texttt{smtpd\_use\_tls} & \texttt{yes} \\
    \texttt{smtpd\_tls\_security\_level} & \texttt{may} \\
    \bottomrule
    \end{tabular}
  \end{tablaAPA}%
}{Elaboraci\'on propia.}
\end{table}

El almacenamiento en formato \texttt{Maildir/} garantiza que cada mensaje
se guarde como un archivo independiente, lo que evita bloqueos al escribir
y simplifica el respaldo. La opci\'on \texttt{noanonymous} impide el env\'io
sin autenticaci\'on.

\capturaAPA{fig:postconf}{Par\'ametros efectivos de Postfix con \texttt{postconf -n}}{captura-postconf.png}{Captura de pantalla de la terminal ejecutando \texttt{postconf -n}, mostrando la configuraci\'on activa del servidor de correo, incluidos los par\'ametros de identidad, redes y SASL.}

\subsection{Configuraci\'on del servidor IMAP/POP3}

Dovecot se configura mediante los archivos del directorio
\texttt{/etc/dovecot/conf.d/}. La configuraci\'on resultante se verificó con
el comando \texttt{doveconf -n}.

\capturaAPA{fig:doveconf}{Verificaci\'on de la configuraci\'on de Dovecot con \texttt{doveconf -n}}{captura-doveconf.png}{Captura de pantalla de la terminal ejecutando \texttt{doveconf -n}, con los par\'ametros activos de Dovecot referidos a protocolos, ubicaci\'on del buz\'on y cifrado.}

\subsection{Pruebas de conectividad y de env\'io}

Desde el equipo cliente se verific\'o la apertura de los puertos del
servicio de correo:

\begin{lstlisting}[style=terminal]
$ for p in 25 110 143 587 993 995; do
    timeout 3 bash -c "echo > /dev/tcp/192.168.0.2/$p" \
      && echo "puerto $p: ABIERTO"
  done
puerto 25: ABIERTO
puerto 587: ABIERTO
\end{lstlisting}

\capturaAPA{fig:puertos-abiertos}{Comprobaci\'on de conectividad a los puertos del servidor de correo}{captura-puertos.png}{Captura de pantalla mostrando la salida del bucle que verifica los puertos 25, 110, 143, 587, 993 y 995 desde la m\'aquina cliente.}

La negociaci\'on \textsc{starttls} se comprob\'o con \texttt{openssl},
que muestra el sujeto del certificado y el c\'odigo de verificaci\'on 18
correspondiente a un certificado auto-firmado:

\begin{lstlisting}[style=terminal]
$ openssl s_client -connect 192.168.0.2:587 -starttls smtp
subject=CN=mail.sudoers.lan, O=TSO-II, C=AR
Verify return code: 18 (self-signed certificate)
\end{lstlisting}

\capturaAPA{fig:starttls}{Negociaci\'on \textsc{starttls} con el servidor SMTP}{captura-starttls.png}{Captura de pantalla de \texttt{openssl s\_client} mostrando el sujeto del certificado y el c\'odigo de verificaci\'on 18, propio de un certificado auto-firmado.}

El env\'io autenticado se valid\'o con \texttt{swaks}, que completa el
intercambio \textsc{smtp} y obtiene la respuesta \texttt{250} de aceptaci\'on
del mensaje:

\begin{lstlisting}[style=terminal]
$ swaks --to joaquin@sudoers.lan \
    --from nicolas@sudoers.lan \
    --server 192.168.0.2 --port 587 \
    --auth LOGIN --auth-user nicolas \
    --header "Subject: Prueba" \
    --body "Hola Joaquin, esto es una prueba."
\end{lstlisting}

\capturaAPA{fig:swaks}{Env\'io de correo autenticado con \texttt{swaks}}{captura-swaks.png}{Captura de pantalla de \texttt{swaks} mostrando el intercambio SMTP, la autenticaci\'on exitosa y el c\'odigo 250 de aceptaci\'on del mensaje.}

La recepci\'on se confirm\'o inspeccionando el buz\'on del destinatario,
donde el mensaje aparece como un archivo individual dentro del directorio
\texttt{new/} del \texttt{Maildir}:

\begin{lstlisting}[style=terminal]
$ ls /home/joaquin/Maildir/new/
1627...joaquin@sudoers.lan
\end{lstlisting}

\capturaAPA{fig:maildir-receptor}{Buz\'on \texttt{Maildir} del destinatario con el mensaje recibido}{captura-maildir.png}{Captura de pantalla listando el contenido de \texttt{/home/joaquin/Maildir/new/}, donde aparece el archivo del correo recibido con el nombre generado por Postfix.}

\subsection{Acceso desde clientes}

El acceso desde un cliente gr\'afico como Thunderbird se realiz\'o con la
configuraci\'on que se detalla en la tabla \ref{tab:config-thunderbird}.

\begin{table}[H]
\centering
\caption{Configuraci\'on del cliente de correo Thunderbird}
\label{tab:config-thunderbird}
\contenidoConFuente{%
  \begin{tablaAPA}
    \begin{tabular}{@{}p{5.4cm}p{8.4cm}@{}}
    \toprule
    \textbf{Par\'ametro} & \textbf{Valor} \\
    \midrule
    Servidor de entrada (IMAP) & \texttt{mail.sudoers.lan} \\
    Puerto de entrada & \texttt{143} \\
    Seguridad de entrada & \textsc{starttls} \\
    Servidor de salida (SMTP) & \texttt{mail.sudoers.lan} \\
    Puerto de salida & \texttt{587} \\
    Seguridad de salida & \textsc{starttls} \\
    Autenticaci\'on & Contrase\~na normal \\
    Usuario & \texttt{joaquin} \\
    \bottomrule
    \end{tabular}
  \end{tablaAPA}%
}{Elaboraci\'on propia.}
\end{table}

\capturaAPA{fig:thunderbird}{Cliente Thunderbird conectado al servidor IMAP con los mensajes recibidos}{captura-thunderbird.png}{Captura de pantalla del cliente Thunderbird mostrando la cuenta \texttt{joaquin@sudoers.lan} conectada y la bandeja de entrada con el mensaje de prueba.}

% ============================================================================
%  BLOQUE 2: SERVIDOR DE ARCHIVOS MULTI-PROTOCOLO
% ============================================================================
\section{Servidor de archivos multi-protocolo}

El servicio de archivos integra tres protocolos --\textsc{smb/cifs} con
Samba, \textsc{nfs} con el m\'odulo del n\'ucleo y \textsc{http} con
Filebrowser-- que operan sobre un mismo directorio compartido
(\texttt{/srv/share}). De este modo, un archivo creado desde un protocolo es
inmediatamente accesible desde los dem\'as, sin duplicar la informaci\'on
almacenada.

\subsection{Configuraci\'on de Samba}

La configuraci\'on reside en \texttt{/etc/smb.conf}. Las directivas
esenciales y su justificaci\'on se resumen en la tabla
\ref{tab:directivas-samba}.

\begin{table}[H]
\centering
\caption{Directivas principales de la configuraci\'on de Samba}
\label{tab:directivas-samba}
\contenidoConFuente{%
  \begin{tablaAPA}
    \begin{tabular}{@{}p{5.4cm}p{8.4cm}@{}}
    \toprule
    \textbf{Directiva} & \textbf{Justificaci\'on} \\
    \midrule
    \texttt{security = user} & Exige autenticaci\'on con usuario y clave registrados \\
    \texttt{map to guest = bad user} & Rechaza usuarios inv\'alidos sin otorgar acceso de invitado \\
    \texttt{guest ok = no} & Proh\'ibe el acceso an\'onimo al recurso \\
    \texttt{valid users} & Lista blanca de los usuarios autorizados \\
    \texttt{create mask / directory mask} & \texttt{0666} y \texttt{0777}: los archivos de un protocolo quedan accesibles a los dem\'as \\
    \bottomrule
    \end{tabular}
  \end{tablaAPA}%
}{Elaboraci\'on propia.}
\end{table}

\begin{lstlisting}[style=terminal]
[compartido]
    path = /srv/share
    browseable = yes
    writable = yes
    guest ok = no
    valid users = nicolas joaquin david
    create mask = 0666
    directory mask = 0777
\end{lstlisting}

El acceso desde la terminal se verific\'o con \texttt{smbclient}, que
autentica al usuario y permite listar y crear contenido en el recurso
compartido:

\capturaAPA{fig:archivos-smbclient}{Conexi\'on al recurso compartido con \texttt{smbclient}}{captura-archivos-smbclient.png}{Captura de pantalla de una sesi\'on \texttt{smbclient} contra el recurso compartido, autenticada con el usuario david y mostrando la lista de archivos y la creaci\'on de un directorio.}

\subsection{Configuraci\'on de NFS}

La exportaci\'n se define en \texttt{/etc/exports}. La opci\'on
\texttt{fsid=0} establece el directorio como ra\'iz virtual del espacio
\textsc{nfsv4}, por lo que el cliente monta la ruta remota \texttt{/} y no
\texttt{/srv/share}:

\begin{lstlisting}[style=terminal]
/srv/share *(rw,sync,no_subtree_check,no_root_squash,fsid=0)

$ showmount -e 192.168.0.2
$ sudo mount -t nfs4 192.168.0.2:/ /mnt/servidor_archivos
\end{lstlisting}

La opci\'on \texttt{sync} prioriza la integridad de los datos al confirmar
en disco antes de responder, mientras que \texttt{no\_subtree\_check}
optimiza el rendimiento en las renombraciones.

\capturaAPA{fig:archivos-nfs}{Verificaci\'on de exportaciones y montaje \textsc{nfs}}{captura-archivos-nfs.png}{Captura de pantalla mostrando la salida de \texttt{showmount -e} y el montaje de \texttt{/} en el punto de acceso del cliente.}

\subsection{Interfaz web de archivos}

Filebrowser expone el mismo directorio \texttt{/srv/share} a trav\'es de
una interfaz web accesible desde cualquier navegador, sin instalar software
adicional en el cliente. La inicializaci\'on de su base de datos y la
creaci\'on de usuarios se realizan al desplegar el servicio:

\begin{lstlisting}[style=terminal]
$ filebrowser config init --database /database/filebrowser.db
$ filebrowser config set --database /database/filebrowser.db \
    --root /srv/share --address 0.0.0.0 --port 80 \
    --auth.method json
\end{lstlisting}

\capturaAPA{fig:archivos-filebrowser}{Interfaz web de Filebrowser autenticada}{captura-archivos-filebrowser.png}{Captura de pantalla del navegador mostrando la pantalla de inicio de sesi\'on y el explorador de archivos sobre el directorio \texttt{/srv/share}.}

% ============================================================================
%  BLOQUE 3: SERVIDOR DHCP
% ============================================================================
\section{Servidor \textsc{dhcp}}

El servicio de asignaci\'on din\'amica de direcciones se implement\'o con
Kea, el servidor \textsc{dhcp} de ISC y sucesor del cl\'asico
\texttt{dhcpd}, seleccionado por su soporte activo, su configuraci\'on
basada en \textsc{json} y su manejo eficiente de las concesiones
\cite{rfc2131, keaDocs}.

\subsection{Arquitectura y segmentaci\'on de direcciones}

El servicio se despleg\'o de manera centralizada sobre el host \texttt{dc1}.
Para que \'unicamente este servidor otorgue direcciones, se desactiv\'o el
servicio \textsc{dhcp} del router, evitando que dos servidores respondieran a
las peticiones \textsc{discover} con ofertas distintas.

La red se dividi\'o en dos segmentos l\'ogicos seg\'un la pol\'itica del
proyecto: las direcciones inferiores a \texttt{192.168.0.110} corresponden a
los equipos de rol, que requieren \textsc{ip} fija, mientras que el rango
superior se reserva para la asignaci\'on din\'amica de los equipos de
usuario. La figura \ref{fig:arquitectura-dhcp} representa el esquema del
servicio.

\figuraAPA{fig:arquitectura-dhcp}{Arquitectura del servicio \textsc{dhcp}: clientes, router y servidor}{%
  \begin{tikzpicture}[
      node distance=1.2cm and 1.6cm,
      ser/.style={draw, rounded corners, minimum width=2.6cm,
                  minimum height=0.9cm, align=center, font=\small},
      cli/.style={draw, minimum width=2.4cm, minimum height=0.8cm,
                  align=center, font=\small},
      flujo/.style={-Stealth, thick}
    ]
    \node[ser] (kea) {Kea DHCP\\\texttt{192.168.0.2}};
    \node[cli, left=of kea, yshift=0.7cm] (c1) {Equipo de usuario\\\texttt{.100+}};
    \node[cli, left=of kea, yshift=-0.7cm] (c2) {Equipo de rol\\\texttt{.50--.99}};
    \node[ser, right=of kea] (router) {Router\\\texttt{192.168.0.1}\\\footnotesize DHCP off};
    \node[ser, below=0.9cm of kea] (dns) {DNS\\\texttt{192.168.0.2}\\\footnotesize mismo host};
    \draw[flujo] (c1.east) -- (kea.west);
    \draw[flujo] (c2.east) -- (kea.west);
    \draw[flujo] (kea.east) -- (router.west);
    \draw[flujo, dashed] (dns.north) -- (kea.south);
  \end{tikzpicture}%
}{Elaboraci\'on propia.}

\subsection{Instalaci\'on y configuraci\'on del servidor}

\begin{lstlisting}[style=terminal]
$ apt update
$ apt install -y kea-dhcp4-server
\end{lstlisting}

La configuraci\'on reside en \texttt{/etc/kea/kea-dhcp4.conf}, en formato
\textsc{json}. El bloque \texttt{subnet4} define el \'ambito, el conjunto de
direcciones y las opciones de red que se entregan al cliente:

\begin{lstlisting}[style=terminal]
"interfaces-config": { "interfaces": ["*"] },
"lease-database": { "type": "memfile", "lfc-interval": 3600 },
"subnet4": [{
  "subnet": "192.168.0.0/24",
  "pools": [{ "pool": "192.168.0.100 - 192.168.0.199" }],
  "option-data": [
    { "name": "domain-name-servers", "data": "192.168.0.2" },
    { "name": "domain-name", "data": "sudoers.lan" },
    { "name": "routers", "data": "192.168.0.1" }
  ]
}]
\end{lstlisting}

El par\'ametro \texttt{lfc-interval} indica cada cu\'antos segundos se
compacta el archivo de concesiones. Las opciones de \texttt{option-data}
garantizan que cada cliente reciba, sin intervenci\'on manual, la m\'ascara
de subred, la puerta de enlace, el servidor \textsc{dns} y el dominio de
b\'usqueda del proyecto.

\capturaAPA{fig:dhcp-conf}{Archivo de configuraci\'on \texttt{kea-dhcp4.conf} con el \'ambito y las opciones de red}{captura-dhcp-conf.png}{Captura de pantalla del archivo de configuraci\'on de Kea mostrando el bloque \texttt{subnet4} con el \'ambito din\'amico, el servidor DNS, el dominio y la puerta de enlace definidos.}

\subsection{Desactivaci\'on del \textsc{dhcp} del router}

El router de la red manten\'ia su propio servicio \textsc{dhcp} activo,
lo que generaba un conflicto: tanto el router como Kea respond\'ian a los
paquetes \textsc{discover}. Para garantizar que el \'unico servidor
\textsc{dhcp} de la red sea el implementado en el proyecto, se accedi\'o a
la interfaz de administraci\'on del router y se deshabilit\'o el servicio en
la \textsc{lan}, manteniendo la configuraci\'on de la puerta de enlace.

\capturaAPA{fig:router-disable}{Interfaz de administraci\'on del router con el servidor \textsc{dhcp} deshabilitado}{captura-router-disable.png}{Captura de pantalla de la configuraci\'on \textsc{lan} del router TP-Link mostrando la opci\'on \textsc{dhcp server} en estado deshabilitado.}

\subsection{Verificaci\'on y pruebas de asignaci\'on}

Antes de probar con clientes reales se valid\'o la sintaxis de la
configuraci\'on y el estado del servicio:

\begin{lstlisting}[style=terminal]
$ kea-dhcp4 -t /etc/kea/kea-dhcp4.conf
Configuration check successful
$ systemctl restart kea-dhcp4-server
$ systemctl status kea-dhcp4-server
\end{lstlisting}

\capturaAPA{fig:estado-kea}{Estado del servicio \texttt{kea-dhcp4-server} en el servidor}{captura-systemctl-kea.png}{Captura de pantalla de la terminal mostrando el estado activo del servicio \texttt{kea-dhcp4-server} con la ruta de la configuraci\'on cargada.}

Desde un equipo cliente se solicit\'o una concesion y se verific\'o la
configuraci\'on recibida:

\begin{lstlisting}[style=terminal]
$ sudo dhclient -v eth0
$ ip a
$ ip route
\end{lstlisting}

\capturaAPA{fig:cliente-ip}{Direcci\'on asignada al cliente tras el arranque de \texttt{dhclient}}{captura-cliente-ip.png}{Captura de pantalla de la salida de \texttt{ip a} en el cliente mostrando la direcci\'on IP dentro del \'ambito din\'amico, la m\'ascara \texttt{/24} y el estado de la interfaz.}

Cada asignaci\'n queda registrada en el archivo de concesiones, que
identifica la direcci\'n entregada, la direcci\'n \textsc{mac} del cliente y
los tiempos de validez:

\capturaAPA{fig:leases}{Concesiones registradas en \texttt{kea-leases4.csv}}{captura-leases.png}{Captura de pantalla del archivo de concessions mostrando las concesiones activas: direcci\'on IP, direcci\'on MAC del cliente, \texttt{hostname} y tiempos de renovaci\'on y expiraci\'on.}

Finalmente se verific\'o la conectividad del cliente hacia la puerta de
enlace y hacia el servidor \textsc{dns}, lo que confirma que la
configuraci\'n entregada por \textsc{dhcp} permite alcanzar los servicios
internos:

\capturaAPA{fig:ping-gateway}{Prueba de conectividad desde el cliente hacia la puerta de enlace}{captura-ping-router-gateway.png}{Captura de pantalla de la terminal del cliente mostrando las respuestas de \texttt{ping} hacia la puerta de enlace \texttt{192.168.0.1} y hacia el servidor del proyecto.}

% ============================================================================
%  BLOQUE 4: SERVIDOR WEB CON PROXY INVERSO
% ============================================================================
\section{Servidor Web con proxy inverso}

El servidor web se implement\'o con Nginx, que cumple un doble rol: aloja
el sitio institucional del proyecto y act\'ua como proxy inverso
centralizado hacia los servicios internos que operan en puertos no
publicados, como el gestor de impresi\'on \textsc{cups}, el explorador de
archivos, el panel de gesti\'on de contenedores y el servidor de monitoreo
\cite{nginxDocs}.

\subsection{Arquitectura del servicio}

Todo el tr\'afico web entrante se protege mediante \textsc{https} con
certificados \textsc{tls} que incluyen nombres alternativos en formato
comodi\'n. La figura \ref{fig:arquitectura-web} ilustra el esquema
implementado.

\figuraAPA{fig:arquitectura-web}{Arquitectura del servidor web: Nginx como proxy inverso y servidor de contenido}{%
  \begin{tikzpicture}[
      node distance=1cm and 1.6cm,
      servicio/.style={draw, rounded corners, minimum width=2.4cm,
                       minimum height=0.8cm, align=center, font=\small},
      cliente/.style={draw, circle, minimum size=0.8cm, font=\small},
      linea/.style={-Stealth, thick}
    ]
    \node[servicio, fill=blue!10] (nginx) {Nginx\\{\footnotesize HTTP/HTTPS}};
    \node[servicio, right=2.2cm of nginx, yshift=2cm] (cups) {CUPS\\{\footnotesize :631}};
    \node[servicio, right=2.2cm of nginx, yshift=0.7cm] (files) {Filebrowser\\{\footnotesize :8082}};
    \node[servicio, right=2.2cm of nginx, yshift=-0.7cm] (portainer) {Portainer\\{\footnotesize :9000}};
    \node[servicio, right=2.2cm of nginx, yshift=-2cm] (zabbix) {Zabbix Web\\{\footnotesize :80}};
    \node[servicio, below=of nginx] (html) {/var/www/\\\footnotesize html};
    \node[cliente, left=of nginx] (browser) {Web};
    \draw[linea] (browser) -- node[above,font=\scriptsize]{443} (nginx);
    \draw[linea] (nginx) -- (cups);
    \draw[linea] (nginx) -- (files);
    \draw[linea] (nginx) -- (portainer);
    \draw[linea] (nginx) -- (zabbix);
    \draw[linea] (nginx) -- (html);
  \end{tikzpicture}%
}{Elaboraci\'on propia.}

La tabla \ref{tab:virtual-hosts} resume los \emph{virtual hosts}
configurados, sus dominios y su destino.

\begin{table}[H]
\centering
\caption{\emph{Virtual hosts} configurados en Nginx}
\label{tab:virtual-hosts}
\contenidoConFuente{%
  \begin{tablaAPA}
    \begin{tabular}{@{}p{4.6cm}p{4cm}p{5.2cm}@{}}
    \toprule
    \textbf{Subdominio} & \textbf{Destino} & \textbf{Funci\'on} \\
    \midrule
    \texttt{web.sudoers.lan} & \texttt{/var/www/html} & Sitio institucional del proyecto \\
    \texttt{print.sudoers.lan} & \texttt{127.0.0.1:631} & Proxy inverso al servicio de impresi\'on \\
    \texttt{files.sudoers.lan} & \texttt{127.0.0.1:8082} & Proxy inverso al gestor de archivos \\
    \texttt{portainer.sudoers.lan} & \texttt{127.0.0.1:9000} & Proxy inverso al panel de gesti\'on \\
    \texttt{zabbix.sudoers.lan} & \texttt{192.168.0.3:80} & Proxy inverso al monitoreo \\
    \texttt{\_} (por defecto) & Redirecci\'on a \textsc{https} & Acceso por IP o dominios no declarados \\
    \bottomrule
    \end{tabular}
  \end{tablaAPA}%
}{Elaboraci\'on propia.}
\end{table}

\subsection{Instalaci\'on y arranque del servicio}

\begin{lstlisting}[style=terminal]
$ sudo apt update
$ sudo apt install -y nginx openssl
$ systemctl enable --now nginx
\end{lstlisting}

El paquete \texttt{nginx} aporta el servidor web de arquitectura orientada a
eventos, adecuado para cargas concurrentes, y \texttt{openssl} permite
generar los certificados \textsc{x.509} requeridos para el protocolo
\textsc{https} \cite{opensslDocs}. Las directivas globales se ubican en
\texttt{/etc/nginx/nginx.conf} y cada \emph{virtual host} en
\texttt{/etc/nginx/sites-available/}.

\capturaAPA{fig:web-arranque}{Arranque del servicio Nginx en el servidor}{captura-web-arranque.png}{Captura de pantalla de la terminal mostrando la instalaci\'on de los paquetes y el arranque del servicio \texttt{nginx} con su estado activo.}

\subsection{Aplicaci\'on de \textsc{tls} y pruebas de acceso}

Tras emitir el certificado con \texttt{openssl} y asociarlo al
\emph{virtual host} se verific\'o la negociaci\'n segura y la publicaci\'on
del sitio a trav\'es del proxy inverso.

\capturaAPA{fig:web-cert}{Certificado \textsc{tls} emitido para el servidor web}{captura-web-cert.png}{Captura de pantalla de la salida de \texttt{openssl} mostrando los datos del certificado generado, incluido el nombre alternativo del dominio.}

\capturaAPA{fig:web-proxy}{Regla de proxy inverso configurada en Nginx}{captura-web-proxy.png}{Captura de pantalla de la configuraci\'n del \emph{virtual host} mostrando las directivas \texttt{proxy\_pass} y \texttt{proxy\_set\_header} que reencaminan la petici\'on al servicio interno.}

\capturaAPA{fig:web-sitio}{Sitio institucional publicado a trav\'es del proxy inverso}{bancosudoers.png}{Captura de pantalla del navegador mostrando el sitio web del proyecto cargado correctamente mediante el subdominio configurado.}

% ============================================================================
%  BLOQUE 5: CONTROLADOR DE DOMINIO SAMBA 4 ACTIVE DIRECTORY
% ============================================================================
\section{Controlador de dominio Samba 4 Active Directory}

El despliegue del dominio constituye el n\'ucleo de identidad de la
infraestructura. Samba 4 integra en un \'unico daemon los protocolos que
forman Active Directory: \textsc{ldap} para el directorio, Kerberos para la
autenticaci\'on y el \textsc{dns} propio del dominio
\cite{sambateam2024, rfc4511, rfc4120}.

\subsection{Preparaci\'on del sistema}

Un controlador de dominio es el punto de referencia de identidad de toda
la red, por lo que su nombre y su direcci\'on \textsc{ip} deben ser estables
y conocidos de antemano. Se fij\'o el nombre \texttt{dc1} y su nombre
completamente cualificado \texttt{dc1.sudoers.lan}:

\begin{lstlisting}[style=terminal]
$ hostnamectl set-hostname dc1
$ hostnamectl status | head -3
  Static hostname: dc1
  Operating System: Debian GNU/Linux 12 (bookworm)
\end{lstlisting}

\capturaAPA{fig:hostname}{Identidad del host \texttt{dc1} configurada en el servidor}{captura-hostname.png}{Captura de pantalla mostrando la salida de \texttt{hostnamectl status} con el nombre de host \texttt{dc1} y la distribuci\'on del sistema operativo.}

Un controlador de dominio no puede tener una direcci\'on variable, ya que
Kerberos y el \textsc{dns} lo referencian constantemente. Se garantiz\'o que
\texttt{dc1} resuelva a \texttt{192.168.0.2/24} en la interfaz cableada y se
registr\'o en \texttt{/etc/hosts}:

\begin{lstlisting}[style=terminal]
$ ip -4 a show eno1
    inet 192.168.0.2/24 brd 192.168.0.255 scope global eno1

$ grep -E 'dc1|sudoers' /etc/hosts
192.168.0.2   dc1.sudoers.lan   dc1
\end{lstlisting}

\capturaAPA{fig:ip-fija}{Direcci\'on \textsc{ip} fija \texttt{192.168.0.2/24} en la interfaz cableada}{captura-ip-fija.png}{Captura de pantalla de la salida de \texttt{ip -4 a show eno1} con la direcci\'on fija del servidor y la entrada de \texttt{/etc/hosts} que asocia \texttt{dc1.sudoers.lan} a esa direcci\'on.}

\subsection{Instalaci\'on de Samba 4 y promoci\'on del servidor}

Los paquetes del dominio se instalaron de forma nativa sobre el host, nunca
dentro de un contenedor: el controlador de Samba es un servicio sensible al
aislamiento de red y sistema de archivos que impone un contenedor
\cite{sambateam2024}.

\begin{lstlisting}[style=terminal]
$ apt update
$ apt install -y samba smbclient winbind krb5-user \
    libpam-winbind libnss-winbind dnsutils
\end{lstlisting}

Junto con Samba se habilitaron los servicios que sustentan el dominio:
el firewall \texttt{nftables}, el servidor de hora \texttt{chrony}
--necesario porque Kerberos rechazaTickets con un desfase de reloj superior
a cinco minutos-- y Docker para los servicios contenedorizados:

\begin{lstlisting}[style=terminal]
$ systemctl is-enabled samba-ad-dc nftables chrony docker
samba-ad-dc.service    enabled
nftables.service       enabled
chrony.service         enabled
docker.service         enabled
\end{lstlisting}

\capturaAPA{fig:servicios-host}{Servicios del host habilitados para el arranque}{captura-servicios-host.png}{Captura de pantalla de la salida de \texttt{systemctl is-enabled} con los servicios \texttt{samba-ad-dc}, \texttt{nftables}, \texttt{chrony} y \texttt{docker} en estado \texttt{enabled}.}

La promoci\'n a controlador de dominio se realiza aprovisionando el dominio
con \texttt{samba-tool}:

\begin{lstlisting}[style=terminal]
$ samba-tool domain provision \
    --use-rfc2307 --realm=SUDOERS.LAN \
    --domain=SUDOERS \
    --server-role=dc \
    --dns-backend=SAMBA_INTERNAL
\end{lstlisting}

La opci\'n \texttt{--use-rfc2307} activa la compatibilidad \textsc{posix}
de los grupos, de modo que la gesti\'on de permisos Unix se integra con la
del dominio.

\subsection{Configuraci\'on del \textsc{dns} del dominio}

El controlador de dominio administra la zona autoritativa interna. La
resoluci\'n se verific\'o con \texttt{dig} tanto para registros de servicio
como para las \textsc{srv} que publican Kerberos esta informaci\'on a
los clientes:

\begin{lstlisting}[style=terminal]
$ dig +short SRV _ldap._tcp.sudoers.lan
$ dig @dc1 sudoers.lan SOA
$ dig @dc1 dc1.sudoers.lan A
\end{lstlisting}

\capturaAPA{fig:zona-dns}{Zona autoritativa \texttt{sudoers.lan} administrada por el controlador de dominio}{captura-zona-dns.png}{Captura de pantalla del contenido de la zona DNS del dominio, con los registros de servicios y de los servidores que la integran.}

\capturaAPA{fig:dig-dc}{Resoluci\'n de nombres contra el controlador de dominio}{captura-dig-dc.png}{Captura de pantalla de la salida de \texttt{dig} con la resoluci\'n de los nombres del dominio contra el servidor \texttt{dc1}.}

\capturaAPA{fig:dig-srv}{Registros \textsc{srv} publicados por el dominio}{captura-dig-srv.png}{Captura de pantalla de la consulta de registros \textsc{srv} del dominio, que publishing las direcciones de los servicios Kerberos y \textsc{ldap}.}

\subsection{Configuraci\'on del servicio \textsc{dhcp} en el dominio}

El servicio \textsc{dhcp} se despliega en un contenedor con la opci\'n
\texttt{network\_mode: host}, condici\'n indispensable porque el protocolo
utiliza la direcci\'n de difusi\'n por \textsc{broadcast} en el puerto
\texttt{67/udp} y no funciona mapeado a trav\'es de un puente de red.

\capturaAPA{fig:kea-ss}{Estado del servicio \textsc{dhcp} del dominio en ejecuci\'on}{captura-kea-ss.png}{Captura de pantalla del estado del servicio de asignaci\'on de direcciones dentro del dominio, confirmar la entrega de direcciones a los clientes.}

\subsection{Administraci\'on de usuarios y grupos}

La gesti\'on de identidades se centraliza en el dominio, lo que evita
credenciales dispersas entre servicios. Las cuentas de administraci\'on se
asignan al grupo \texttt{admins} y las de prueba al grupo de oficina:

\begin{lstlisting}[style=terminal]
$ samba-tool user create joaquin --given-name="Joaquin" \
    --random-password
$ samba-tool group add admins
$ samba-tool group add users
$ samba-tool group add machines
\end{lstlisting}

\capturaAPA{fig:usuarios}{Usuarios creados en el dominio}{captura-usuarios.png}{Captura de pantalla de la salida de \texttt{samba-tool user list} mostrando las cuentas registradas en el dominio.}

\capturaAPA{fig:grupos}{Grupos de seguridad del dominio}{captura-grupos.png}{Captura de pantalla de la salida de \texttt{samba-tool group list} con los grupos definidos, incluidos los de administraci\'on y usuarios.}

\capturaAPA{fig:user-show}{Detalle de una cuenta de usuario del dominio}{captura-user-show.png}{Captura de pantalla de la informaci\'n detallada de una cuenta de usuario, incluyendo su pertenencia a los grupos del dominio.}

\subsection{Uni\'on de clientes y verificaci\'on del dominio}

La uni\'on de un cliente al dominio se verifica obteniendo un ticket de
Kerberos con \texttt{kinit} y list\'andolo con \texttt{klist}:

\begin{lstlisting}[style=terminal]
$ kinit joaquin
Password for joaquin@SUDOERS.LAN:
$ klist
Ticket cache: FILE:/tmp/krb5cc_0
Default principal: joaquin@SUDOERS.LAN
\end{lstlisting}

\capturaAPA{fig:kinit-cliente}{Obtenci\'on de ticket de Kerberos en el cliente}{captura-kinit-cliente.png}{Captura de pantalla de la terminal del cliente ejecutando \texttt{kinit} para obtener un ticket de Kerberos del dominio.}

\capturaAPA{fig:klist}{Listado de tickets de Kerberos del cliente}{captura-kinit-klist.png}{Captura de pantalla de la salida de \texttt{klist} mostrando la cach\'e de tickets y el principal por defecto del usuario autenticado.}

La resoluci\'n de grupos en el lado del cliente se verific\'o con
\texttt{getent}, que confirma la integraci\'on de \texttt{winbind} con la
resoluci\'on de nombres:

\capturaAPA{fig:getent}{Resoluci\'on de grupos del dominio en el cliente}{captura-getent.png}{Captura de pantalla de la salida de \texttt{getent group} mostrando los grupos del dominio resueltos correctamente en el cliente.}

El acceso a los recursos compartidos se verific\'o con \texttt{smbclient}
usando las credenciales del dominio:

\capturaAPA{fig:smbclient-ad}{Acceso a recursos compartidos con credenciales del dominio}{captura-smbclient.png}{Captura de pantalla de una sesi\'on \texttt{smbclient} autenticada con las credenciales del dominio, mostrando el acceso al recurso compartido.}

La informaci\'n general del dominio se consult\'o con las utilidades
propias de Samba:

\capturaAPA{fig:domain-info}{Informaci\'on del dominio y de la lista de equipos}{captura-domain-info.png}{Captura de pantalla de la salida de \texttt{samba-tool domain info} con los datos del dominio y su relaci\'on de confianza.}

\capturaAPA{fig:computer-list}{Equipos unidos al dominio}{captura-computer-list.png}{Captura de pantalla de la salida de \texttt{samba-tool computer list} mostrando los equipos que han unido al dominio.}

\subsection{Puertos publicados por el dominio}

El controlador de dominio expone un conjunto definido de puertos, cuya
disponibilidad se verific\'o desde el cliente:

\capturaAPA{fig:puertos-ad}{Verificaci\'on de los puertos de servicio del dominio}{captura-puertos-ad.png}{Captura de pantalla de la terminal del cliente mostrando los puertos del controlador de dominio que responden: DNS, Kerberos, \textsc{ldap} y SMB.}

% ============================================================================
%  BLOQUE 6: DNS DEL DOMINIO Y SERVIDOR WEB INTEGRADO (a8)
% ============================================================================
\section{DNS del dominio y servidor Web integrado}

La \'ultima etapa del proyecto consolid\'o en un \'unico dise\~no los dos
servicios que articulan toda la infraestructura: el \textsc{dns} autoritativo
del dominio \texttt{sudoers.lan}, servido por el propio controlador de
dominio, y el servidor web que publica los servicios de la intranet por
nombre. Este bloque describe la soluci\'n final, donde el \textsc{dns} y la
puerta web convergen sobre una \'unica direcci\'n de entrada \cite{rfc1034}.

\subsection{Arquitectura conjunta \textsc{dns} + \textsc{web}}

Como todos los servicios se ejecutan en el \'unico servidor
(\texttt{192.168.0.2}), el \textsc{dns} del dominio y la puerta web comparten
esa misma direcci\'n. La separaci\'on entre ambos es \textbf{l\'ogica} y no
f\'isica: el \textsc{dns} se atiende de forma nativa desde el daemon del
controlador de dominio, mientras que Nginx corre en un contenedor que
publica los puertos 80 y 443 \cite{tanenbaum2017}.

La tabla \ref{tab:mapa-nombres} resume los nombres gestionados por el
\textsc{dns} del dominio y el servicio al que dirige cada uno el proxy web.

\begin{table}[H]
\centering
\caption{Mapa de nombres del dominio y su backend en el proxy}
\label{tab:mapa-nombres}
\contenidoConFuente{%
  \begin{tablaAPA}
    \begin{tabular}{@{}p{5.2cm}p{4.4cm}p{4.2cm}@{}}
    \toprule
    \textbf{Nombre} & \textbf{Backend} & \textbf{Registro} \\
    \midrule
    \texttt{sudoers.lan}, \texttt{web}, \texttt{www} & Portal est\'atico (Nginx) & A \\
    \texttt{dc1.sudoers.lan} & Controlador de dominio (DNS, SMB) & A \\
    \texttt{print.sudoers.lan} & \texttt{tso-print:631} (CUPS) & A \\
    \texttt{portainer.sudoers.lan} & \texttt{tso-portainer:9000} & A \\
    \texttt{joaquin.sudoers.lan} & P\'agina del equipo (est\'atico) & A \\
    \texttt{david.sudoers.lan} & \texttt{192.168.0.51:8080} & A \\
    \texttt{nicolas.sudoers.lan} & \texttt{192.168.0.52:8080} & A \\
    \texttt{\_ldap.\_tcp.dc.\_msdcs} & \texttt{dc1}:389 (LDAP) & SRV \\
    \bottomrule
    \end{tabular}
  \end{tablaAPA}%
}{Elaboraci\'on propia.}
\end{table}

\subsection{El \textsc{dns} autoritativo servido por Samba}

Un dominio de Active Directory depende por completo del \textsc{dns}: los
clientes lo consultan para descubrir d\'onde reside cada servicio, tanto
mediante nombres comunes como mediante registros de descubrimiento
\textsc{srv} \cite{rfc2782, msft_activedirectory}. Por esta raz\'on, durante
el aprovisionamiento se indic\'o expl\'icitamente que el propio Samba
administrase la zona mediante el backend \textsc{samba\_internal}:

\begin{lstlisting}[style=terminal]
$ samba-tool domain provision \
    --realm=SUDOERS.LAN --domain=SUDOERS \
    --server-role=dc --dns-backend=SAMBA_INTERNAL \
    --use-rfc2307 --host-name=dc1 --host-ip=192.168.0.2 \
    --option="dns forwarder = 192.168.0.1"
\end{lstlisting}

Con este backend los registros se almacenan dentro del directorio y se
actualizan de forma din\'amica, sin necesidad de archivos de zona
manuales. La configuraci\'n resultante queda registrada en
\texttt{/etc/samba/smb.conf}:

\begin{lstlisting}[style=terminal]
$ grep -E 'server role|dns|realm' /etc/samba/smb.conf
[global]
    server role = active directory domain controller
    realm = SUDOERS.LAN
    dns forwarder = 192.168.0.1
    dns update command = /usr/sbin/samba_dnsupdate
\end{lstlisting}

\capturaAPA{fig:smb-dns}{Zona \texttt{sudoers.lan} y rol del controlador de dominio en \texttt{smb.conf}}{captura-smb-dns.png}{Captura de pantalla mostrando la salida del \texttt{grep} sobre \texttt{/etc/samba/smb.conf} con el rol \texttt{active directory domain controller}, el \textit{realm} \texttt{SUDOERS.LAN} y la directiva \texttt{dns forwarder = 192.168.0.1}.}

\subsection{Registros de la zona y de descubrimiento}

El contenido \'integral de la zona se consulta con \texttt{samba-tool dns},
que devuelve los registros \textsc{a} de todos los servicios publicados:

\begin{lstlisting}[style=terminal]
$ kinit administrator
$ samba-tool dns query dc1.sudoers.lan sudoers.lan @ ALL -k yes
Name=, Records=8, Children=0
  DC@SUDOERS.LAN: type=A, value=192.168.0.2
  web@SUDOERS.LAN: type=A, value=192.168.0.2
  print@SUDOERS.LAN: type=A, value=192.168.0.2
  portainer@SUDOERS.LAN: type=A, value=192.168.0.2
\end{lstlisting}

\capturaAPA{fig:zona-dns-a8}{Zona DNS completa del dominio consultada con \texttt{samba-tool dns}}{captura-zona-dns-a8.png}{Captura de pantalla mostrando la salida de \texttt{samba-tool dns query dc1.sudoers.lan sudoers.lan @ ALL} con los registros A de \texttt{dc1}, \texttt{web}, \texttt{print}, \texttt{portainer} y de los administradores, todos hacia \texttt{192.168.0.2}.}

La resoluci\'n de un nombre individual se verifica con \texttt{dig}, que
responde con la direcci\'n del proxy:

\begin{lstlisting}[style=terminal]
$ dig web.sudoers.lan
;; ANSWER SECTION:
web.sudoers.lan. 900 IN A 192.168.0.2
\end{lstlisting}

\capturaAPA{fig:dig-web}{Resoluci\'on de \texttt{web.sudoers.lan} por el \textsc{dns} del dominio}{captura-dig-web.png}{Captura de pantalla mostrando la salida de \texttt{dig web.sudoers.lan} con la respuesta \texttt{A 192.168.0.2}.}

Los registros \textsc{srv} de descubrimiento fueron generados
autom\'aticamente durante el aprovisionamiento y son los que permiten a los
clientes localizar el \textsc{ldap} del dominio:

\begin{lstlisting}[style=terminal]
$ dig _ldap._tcp.dc._msdcs.sudoers.lan SRV
;; ANSWER SECTION:
_ldap._tcp.dc._msdcs.sudoers.lan. 900 IN SRV 0 100 389 dc1.sudoers.lan.
\end{lstlisting}

\capturaAPA{fig:dig-srv-a8}{Registros \textsc{srv} de descubrimiento generados por el controlador de dominio}{captura-dig-srv-a8.png}{Captura de pantalla mostrando la salida de \texttt{dig \_ldap.\_tcp.dc.\_msdcs.sudoers.lan SRV} con el registro SRV apuntando al puerto 389 de \texttt{dc1.sudoers.lan}.}

\subsection{Registros de los equipos de administraci\'on}

Los nombres de los miembros del equipo tambi\'en resuelven hacia el proxy y
no hacia la direcci\'n de sus m\'aquinas. Gracias a esta indirecci\'on, un
backend puede cambiar de m\'aquina o de puerto sin que los nombres ni los
clientes se modifiquen. Los registros se crean con un script idempotente
que puede ejecutarse varias veces sin duplicarlos:

\begin{lstlisting}[style=terminal]
$ sudo kinit administrator
$ sudo bash dc/dns-records.sh
==> Servicios web alojados en maquinas de admins (via proxy) ...
  ok    david.sudoers.lan    ->  192.168.0.2  (ya existe)
  ok    nicolas.sudoers.lan  ->  192.168.0.2  (ya existe)
\end{lstlisting}

\capturaAPA{fig:dns-records}{Creaci\'on idempotente de los registros de los administradores}{captura-dns-records.png}{Captura de pantalla mostrando la salida de \texttt{sudo bash dc/dns-records.sh} con los registros \texttt{david} y \texttt{nicolas} apuntando a \texttt{192.168.0.2}.}

\subsection{Reenv\'io hacia Internet}

Para que los clientes resuelvan tambi\'en nombres externos, el controlador
act\'ua como \textit{forwarder} hacia el router de la red. De esta manera una
\'unica consulta al \textsc{dns} del dominio cubre tanto los nombres internos
como los de Internet, sin configurar un segundo servidor de nombres:

\begin{lstlisting}[style=terminal]
$ dig google.com @192.168.0.2 +short
142.250.78.78          # resolucion externa via el forwarder 192.168.0.1
\end{lstlisting}

\capturaAPA{fig:dig-externo}{Resoluci\'on de nombres externos a trav\'es del \textit{forwarder}}{captura-dig-externo.png}{Captura de pantalla mostrando la salida de \texttt{dig google.com @192.168.0.2} con una direcci\'on IP p\'ublica resuelta.}

\subsection{Integraci\'on con el servicio \textsc{dhcp}}

Para que una estaci\'n reci\'en conectada quede operativa sin intervenci\'on
manual, el servidor \textsc{dhcp} entrega junto con la direcci\'n las opciones
de red que apuntan al \textsc{dns} del dominio: la direcci\'n
\texttt{192.168.0.2}, el dominio \texttt{sudoers.lan}, la puerta de enlace
\texttt{192.168.0.1} y el servidor de hora. Toda m\'aquina de la red queda
as\'i capaz de resolver los nombres internos de forma transparente.

\subsection{Nginx como contenedor y proxy inverso}

El servidor web se despliega como servicio del \texttt{docker-compose} del
proyecto, a diferencia del controlador de dominio que corre de forma
nativa. La definici\'n del servicio establece los puertos publicados y los
vol\'umenes con el portal y los certificados:

\begin{lstlisting}[style=terminal]
web:
  build: ./services/web          # nginx:1.27-alpine
  container_name: tso-web
  ports:
    - "80:80"                    # redirige a HTTPS
    - "443:443"                  # HTTPS (termina TLS)
  volumes:
    - ./services/web/public:/usr/share/nginx/html:ro
    - ./services/web/tls:/etc/nginx/tls:ro      # CA + cert wildcard
  restart: always
\end{lstlisting}

\capturaAPA{fig:compose-web}{Servicio \texttt{web} definido en el \texttt{docker-compose.yml}}{captura-compose-web.png}{Captura de pantalla mostrando el bloque \texttt{web:} del \texttt{docker-compose.yml} con los puertos 80/443, los vol\'umenes del portal y de \textsc{tls}, y la directiva \texttt{restart: always}.}

Nginx enruta cada petici\'on por nombre y, adem\'as, redirige todo el
tr\'afico en claro hacia el canal cifrado:

\begin{lstlisting}[style=terminal]
server {
    listen 80;                                 # HTTP: solo redirige
    server_name sudoers.lan web.sudoers.lan www.sudoers.lan;
    return 301 https://$host$request_uri;
}

server {
    listen 443 ssl;                            # HTTPS: portal
    http2 on;
    server_name sudoers.lan web.sudoers.lan www.sudoers.lan;
    ssl_certificate     /etc/nginx/tls/server.crt;
    ssl_certificate_key /etc/nginx/tls/server.key;
    root  /usr/share/nginx/html;
    index index.html;
}
\end{lstlisting}

\capturaAPA{fig:nginx-conf}{Virtual host del portal y redirecci\'on de \textsc{http} a \textsc{https}}{captura-nginx-conf.png}{Captura de pantalla mostrando el contenido de la configuraci\'n de Nginx con el bloque \texttt{server} del portal y la directiva \texttt{return 301 https://\$host\$request\_uri}.}

\subsection{\textsc{https} con autoridad de certificaci\'on interna}

El portal se publica mediante \textsc{https} con un certificado
\textit{wildcard} que cubre todos los subdominios, firmado por una
autoridad de certificaci\'on interna generada con el script
\texttt{services/web/tls-gen.sh}:

\begin{lstlisting}[style=terminal]
$ sudo bash services/web/tls-gen.sh
=> Generando CA interna ...
=> Generando certificado para *.sudoers.lan ...
============================================
  CA:       ca.crt (instalar en los clientes para confiar)
  Server:   server.crt / server.key  (*.sudoers.lan, 825 dias)
============================================
$ docker compose up -d --build web
\end{lstlisting}

\capturaAPA{fig:tls-gen}{Generaci\'on de la autoridad interna y del certificado \textit{wildcard}}{captura-tls-gen.png}{Captura de pantalla mostrando la salida de \texttt{sudo bash services/web/tls-gen.sh} con la creaci\'on de la autoridad de certificaci\'on y del certificado \texttt{*.sudoers.lan}.}

El acceso al portal se verifica con \texttt{curl}, que devuelve el c\'odigo
de respuesta correcto:

\begin{lstlisting}[style=terminal]
$ curl -skI https://web.sudoers.lan
HTTP/1.1 200 OK
Server: nginx
\end{lstlisting}

\capturaAPA{fig:curl-web}{Acceso \textsc{https} al portal servido por Nginx}{captura-curl-web.png}{Captura de pantalla mostrando la salida de \texttt{curl -skI https://web.sudoers.lan} con el c\'odigo \texttt{200 OK} y el encabezado \texttt{Server: nginx}.}

\subsection{Proxy hacia los servicios de los administradores}

Los servicios alojados en los equipos de los administradores se exponen
tambi\'en a trav\'es del proxy. Sus nombres resuelven hacia
\texttt{192.168.0.2} y Nginx alcanza la m\'aquina correspondiente por la red
plana, de modo que el \'unico puerto visible desde la \textsc{lan} es el
443:

\begin{lstlisting}[style=terminal]
server {
    listen 443 ssl;
    server_name david.sudoers.lan;
    location / {
        resolver 127.0.0.11 valid=10s ipv6=off;
        set $david_upstream http://192.168.0.51:8080;
        proxy_pass $david_upstream;
        proxy_connect_timeout 5s;
    }
}
\end{lstlisting}

\begin{lstlisting}[style=terminal]
$ dig david.sudoers.lan +short          # el cliente resuelve el proxy
192.168.0.2
$ curl -sk https://david.sudoers.lan    # y el proxy responde la maquina
... (contenido del servicio de David) ...
\end{lstlisting}

\capturaAPA{fig:curl-david}{Servicio de la m\'aquina de David publicado a trav\'es del proxy}{captura-curl-david.png}{Captura de pantalla mostrando la resoluci\'on \texttt{dig david.sudoers.lan +short} con \texttt{192.168.0.2} y la respuesta de \texttt{curl -sk https://david.sudoers.lan} con el contenido del servicio del administrador.}

\subsection{Pruebas de integraci\'on desde un cliente}

La validaci\'on final se realiz\'o desde un equipo de la \textsc{lan}, que
recibi\'o su direcci\'n y su servidor \textsc{dns} por \textsc{dhcp}, y
confirm\'o que tanto la resoluci\'o de nombres como el acceso web funcionan
sin configuraci\'on manual:

\begin{lstlisting}[style=terminal]
$ dig web.sudoers.lan +short             # desde el cliente
192.168.0.2
$ curl -sk https://web.sudoers.lan | head -5
<!DOCTYPE html>
<html lang="es">
... (portal de la intranet) ...
\end{lstlisting}

\capturaAPA{fig:cliente-dig}{Resoluci\'on de nombres desde un cliente de la red}{captura-cliente-dig.png}{Captura de pantalla mostrando la salida de \texttt{dig web.sudoers.lan} ejecutada en un cliente, respondiendo \texttt{192.168.0.2}.}

\capturaAPA{fig:cliente-web}{Acceso \textsc{https} al portal desde el navegador del cliente}{captura-cliente-web.png}{Captura de pantalla del navegador de un cliente mostrando el portal de la intranet cargado por \texttt{https://web.sudoers.lan} con el candado de confianza de la autoridad interna.}

La misma validaci\'on se repiti\'o para el portal, la consola de impresi\'on,
el panel de gesti\'on de contenedores y las p\'aginas de los miembros, lo que
confirm\'o que el \textsc{dns} y el proxy web operan como un \'unico sistema
de nombres y de publicaci\'on.

% ============================================================================
%  BLOQUE 7: MEDIDAS DE SECURIZACI\'ON
% ============================================================================
\section{Medidas de seguridad aplicadas}

La seguridad de la infraestructura se abord\'o de forma transversal,
aplicando controles espec\'ificos en cada capa de la arquitectura descrita
en el cap\'itulo anterior \cite{stewart2017}.

\subsection{Seguridad perimetral y control de acceso}

El firewall \texttt{nftables} opera con pol\'itica \textsc{drop} por
defecto, de modo que todo el tr\'afico no permitido se descarta
expl\'icitamente. El acceso al puerto \texttt{22/tcp} se restringe a los
rango de direcciones reservadas para los administradores de TI, de modo
que los puestos operativos no pueden alcanzar la consola del servidor:

\begin{lstlisting}[style=terminal]
# Politica por defecto: descartar todo
$ nft add table inet filter
$ nft add chain inet filter input '{ type filter hook input priority 0; policy drop; }'

# Permitir administracion solo desde el rango de TI
$ nft add rule inet filter input ip saddr 192.168.0.50-192.168.0.52 \
    tcp dport 22 ct state new accept
\end{lstlisting}

\subsection{Aislamiento de servicios internos}

Los servicios que no deben exponerse a la red se mantienen en puertos
internos y se publican \'unicamente a trav\'es del proxy inverso Nginx. De
esta forma, la base de datos \textsc{postgresql} y el panel de gesti\'on de
contenedores no son alcanzables de forma directa desde la \textsc{lan}.

\subsection{Autenticaci\'on y control de acceso por rol}

La gesti\'on centralizada de identidades en el dominio permite asignar
permisos a los recursos compartidos por grupo de usuarios, en lugar de
otorgar acceso individual a cada estaci\'on. La autorizaci\'on del correo
se restricti\'a a los usuarios autenticados mediante \textsc{sasl}, y el
env\'io an\'onimo queda deshabilitado con la opci\'on
\texttt{noanonymous}.

\subsection{Cifrado del tr\'afico}

Las comunicaciones sensibles se protegen con \textsc{tls} mediante
certificados emitidos para los dominios internos, lo que impide el
tr\'ansito de credenciales en texto claro tanto en el env\'io de correo
como en el acceso web.

% ============================================================================
%  BLOQUE 7: PRUEBAS Y VALIDACI\'ON
% ============================================================================
\section{Pruebas funcionales y validaci\'on de la infraestructura}

La validaci\'on final se realiz\'o comprobando la resolutividad de los
nombres del dominio, la obtenci\'on de una direcci\'on v\'alida por
\textsc{dhcp} y el acceso a los servicios publicados a trav\'es del proxy
inverso. Las pruebas ejecutadas desde el cliente fueron:

\begin{lstlisting}[style=terminal]
$ ping -c 4 192.168.0.2
$ dig @192.168.0.2 web.sudoers.lan
$ curl -I https://web.sudoers.lan
$ kinit joaquin
\end{lstlisting}

La obtenci\'n de la direcci\'n por \textsc{dhcp} se valid\'o mediante la
salida de \texttt{dhclient}, en la que el cliente renueva la concesion
previa y recibe una direcci\'n dentro del \'ambito din\'amico, junto con la
m\'ascara, la puerta de enlace y el servidor \textsc{dns} del proyecto.
Esta renovaci\'n result\'o especialmente significativa porque la direcci\'n
anterior proven\'ia del servicio \textsc{dhcp} del router, y su
sustituci\'n confirm\'o que el servidor del dominio asume \'unicamente la
totalidad de las concessions de la red.

Los resultados obtenidos confirman que los tres servicios operan de forma
conjunta: el cliente obtiene su direcci\'n, resuelve los nombres del dominio
contra el controlador de dominio y accede a los servicios publicados a
trav\'es del proxy inverso, validando la integraci\'on integral de la
infraestructura.
